% Part: sets-functions-relations
% Chapter: relations
% Section: equivalence-relations

\documentclass[../../../include/open-logic-section]{subfiles}

\begin{document}

\olfileid[tr]{sfr}{rel}{eqv}

\olsection{Denklik Bağıntıları}

Bir küme üzerindeki özdeşlik bağıntısı yansımalı, simetrik ve geçişlidir. Bu üç
özelliğin hepsini taşıyan $R$ bağıntılarına çok sık rastlanır.

\begin{defn}[Denklik bağıntısı]
Yansımalı, simetrik ve geçişli bir $R \subseteq A^2$ bağıntısına \emph{denklik
bağıntısı} denir. $A$'nın $x$ ve $y$ elemanlarına, $Rxy$ ise
\emph{$R$-denk} denir.
\end{defn}

Denklik bağıntıları \emph{denklik sınıfı} kavramını doğurur. Bir denklik
bağıntısı tanım kümesini birbirinden ayrı parçalara ayırır. Her parçadaki bütün
nesneler birbiriyle bağıntılıdır; farklı parçalardaki hiçbir nesne birbiriyle
bağıntılı değildir. Kimi zaman doğrudan bu parçalardan söz etmek yararlı olur.
Bu amaçla şu tanımı getiriyoruz:

\begin{defn}\ollabel{def:equivalenceclass}
$R \subseteq A^2$ bir denklik bağıntısı olsun. Her $x \in A$ için $x$'in
$A$'daki \emph{denklik sınıfı},
$\equivrep{x}{R}=\Setabs{y \in A}{Rxy}$ kümesidir. $A$'nın $R$'ye göre
\emph{bölüm kümesi},
$\equivclass{A}{R}=\Setabs{\equivrep{x}{R}}{x \in A}$; yani bu denklik
sınıflarının kümesidir.
\end{defn}

Bir sonraki sonuç, denklik sınıflarının gerçekten de $A$'nın parçaları
olduğunu ispatlayarak denklik sınıfı tanımının yerindeliğini gösterir:

\begin{prop}
$R \subseteq A^2$ bir denklik bağıntısıysa, $Rxy$ ancak ve ancak
$\equivrep{x}{R}=\equivrep{y}{R}$ olur.
\end{prop}

\begin{proof}
Soldan sağa yön için $Rxy$ olduğunu varsayalım ve
$z \in \equivrep{x}{R}$ olsun. O hâlde tanım gereği $Rxz$'dir. $R$ bir denklik
bağıntısı olduğundan $Ryz$ olur. (Ayrıntılı yazarsak: $Rxy$ ve $R$'nin
simetrikliği $Ryx$'i; $Rxz$ ve $R$'nin geçişliliği de $Ryz$'yi verir.)
Dolayısıyla $z \in \equivrep{y}{R}$. Genelleştirirsek
$\equivrep{x}{R} \subseteq \equivrep{y}{R}$ elde ederiz. Tam olarak aynı
biçimde $\equivrep{y}{R} \subseteq \equivrep{x}{R}$ olur. Kaplamsallık
uyarınca $\equivrep{x}{R}=\equivrep{y}{R}$ sonucuna varırız.

Sağdan sola yön için $\equivrep{x}{R}=\equivrep{y}{R}$ olduğunu varsayalım.
$R$ yansımalı olduğundan $Ryy$ ve dolayısıyla
$y \in \equivrep{y}{R}$ olur. $\equivrep{x}{R}=\equivrep{y}{R}$ varsayımına
göre $y \in \equivrep{x}{R}$ da geçerlidir. Öyleyse $Rxy$'dir.
\end{proof}

\begin{ex}
Denklik bağıntılarının güzel bir örneği modüler aritmetikten gelir. Her
$a$, $b$ ve $n \in \PosInt$ için, $a \equiv_n b$ ancak ve ancak $a$'nın
~$n$'ye bölümünden kalanla $b$'nin~$n$'ye bölümünden kalan aynı olduğunda
geçerli olsun. (Biraz daha simgesel olarak: $a \equiv_n b$ ancak ve ancak
$a-b=kn$ olacak biçimde bir $k \in \Int$ bulunduğunda geçerli olsun.) Bu
durumda her~$n$ için $\equiv_n$ bir denklik
bağıntısıdır. Ayrıca $\equiv_n$'nin oluşturduğu tam olarak $n$ ayrı denklik
sınıfı vardır; yani $\equivclass{\Nat}{\equiv_n}$ kümesinin $n$ elemanı vardır.
Bunlar şunlardır: $n$'ye kalansız bölünebilen sayıların kümesi, yani
$\equivrep{0}{\equiv_n}$; $n$'ye bölümünden kalan~$1$ olan sayıların kümesi,
yani $\equivrep{1}{\equiv_n}$; \ldots; ve $n$'ye bölümünden kalan~$n-1$ olan
sayıların kümesi, yani~$\equivrep{n-1}{\equiv_n}$.
\end{ex}

\begin{prob}
Her $n \in \PosInt$ için $\equiv_n$'nin bir denklik bağıntısı olduğunu ve
$\equivclass{\Nat}{\equiv_n}$ kümesinin tam olarak $n$ elemanı olduğunu
gösterin.
\end{prob}

\end{document}
